% Report Module 5 – Cyber Terrorism and Information Warfare
% CS 475: Introduction to Computer Security
% Mária Nyolcas
% March 19th 2026

\documentclass[11pt,a4paper]{article}

\usepackage[margin=2.5cm]{geometry}
\usepackage{setspace}
\usepackage{booktabs}
\usepackage{tabularx}
\usepackage{array}
\usepackage{parskip}
\usepackage{titlesec}
\usepackage{hyperref}
\usepackage{enumitem}
\usepackage[protrusion=true,expansion=false]{microtype}
\usepackage{fancyhdr}
\usepackage{graphicx}
\usepackage{xcolor}

\setstretch{1.15}

\hypersetup{
    colorlinks=true,
    linkcolor=black,
    urlcolor=blue,
    citecolor=black,
}

\titleformat{\section}{\large\bfseries}{}{0em}{}[\titlerule]
\titleformat{\subsection}{\normalsize\bfseries}{}{0em}{}

\pagestyle{fancy}
\fancyhf{}
\lhead{\small CS 475 – Mária Nyolcas}
\rhead{\small Week 5: Cyber Terrorism \& Information Warfare}
\cfoot{\thepage}
\renewcommand{\headrulewidth}{0.4pt}

\begin{document}

\begin{center}
  {\Large\bfseries Report Module 5}\\[4pt]
  {\large Cyber Terrorism and Information Warfare}\\[6pt]
  {\normalsize CS 475: Introduction to Computer Security}\\[2pt]
  {\normalsize Mária Nyolcas \quad|\quad Prof: Asen Nikolov Grozdanov}\\[2pt]
  {\normalsize March 2026}
\end{center}

\vspace{4pt}
\noindent\rule{\linewidth}{0.6pt}
\vspace{2pt}

\noindent\textbf{Abstract.}
This report examines cyber terrorism and information warfare as distinct but
overlapping threat categories, combining theoretical analysis with three practical
exercises: threat actor mapping using MITRE ATT\&CK, an anatomy of a documented
disinformation campaign (Operation Secondary Infektion), and the design of a
fictional attack scenario against a hospital network with paired countermeasures.
The report connects these topics to vulnerability and scanning work from Week~4
and ends with a forward link to industrial espionage.


\section{Introduction}

Cyber terrorism and information warfare are two of the most debated concepts in
contemporary security. The confusion is partly terminological: the word
``terrorism'' carries legal and political weight that shapes how incidents are
classified, investigated, and attributed. A working definition that cuts through
the noise comes from RAND Corporation: cyber terrorism consists of
``premeditated, politically motivated attacks against information, computer
systems, computer programs, and data which result in violence against
non-combatant targets by sub-national groups or clandestine agents.'' The key
ingredients are intent (political), target (civilian or non-combatant
infrastructure), and effect (fear, coercion, or physical disruption through
digital means).

This distinguishes cyber terrorism from two closely related categories.
\textbf{Cybercrime} is financially motivated: ransomware gangs, payment card
fraud, and cryptocurrency theft are canonical examples. \textbf{Hacktivism} is
ideologically motivated but typically non-violent: website defacement and
DDoS protests are disruptive but rarely cause lasting physical harm. The
boundary between the three is blurry in practice --- a ransomware attack on a
hospital may simultaneously constitute cybercrime (extortion), potential
terrorism (endangering patients), and, if state-sponsored, an act of warfare.

\textbf{Information Warfare (IW)} predates the internet by centuries; propaganda
leaflets dropped from aircraft in WWI are textbook examples. What digital
networks changed is scale, speed, and deniability. Modern IW doctrine covers six
sub-domains: psychological operations (PSYOP), electronic warfare, cyber
operations, deception, OSINT exploitation, and censorship or narrative control.
This module sits at the intersection of all six, making it the most cross-cutting
topic in the course so far.

For a computer science student, the relevance is concrete: the same scanning
and exploitation techniques studied in Week~4 (Nmap, OpenVAS, CVE analysis)
are the operational foundation for the cyber operations sub-domain of IW. The
distinction is motivation and scale, not method.


\section{Theoretical Framework}

\subsection{Actor Taxonomy}

Understanding who conducts cyber terrorism and IW is as important as understanding
the technical methods, because the actor type determines capability, persistence,
and the likelihood of attribution. Four categories matter for threat modelling:

\begin{itemize}[nosep]
  \item \textbf{Nation-state APTs} operate with virtually unlimited resources,
        long planning horizons, and geopolitical objectives. MITRE ATT\&CK documents
        dozens of these groups with their full TTP profiles.
  \item \textbf{Non-state terrorist organisations} increasingly use cyber
        capabilities for recruitment, propaganda distribution, and cryptocurrency
        fundraising --- less often for direct disruption.
  \item \textbf{Hacktivist collectives} are loosely organised and ideologically
        driven; their campaigns often blur into IW territory when they amplify
        disinformation alongside DDoS attacks.
  \item \textbf{Insider threats} are relevant because, as Week~2's access-control
        analysis showed, overly permissive role assignments create the conditions
        that insiders (and APTs using stolen credentials) exploit.
\end{itemize}

\subsection{Landmark Case Studies}

\textbf{Estonia 2007.} A wave of DDoS attacks paralysed Estonian government,
banking, and media websites following a political dispute with Russia over the
relocation of a Soviet-era war memorial. This is the first large-scale cyber
attack against a nation-state's digital infrastructure and directly accelerated
the creation of NATO CCDCOE in Tallinn. The key takeaway is that even
technically unsophisticated attacks (volumetric DDoS) can cause strategic-level
disruption if the target lacks redundancy and has not planned for the scenario.

\textbf{NotPetya (2017).} Initially disguised as ransomware, NotPetya was a
wiper --- malware designed to permanently destroy data with no recovery
mechanism. Attributed to the Sandworm APT group (Russian GRU Unit 74455), it
caused an estimated \$10 billion in global damages by spreading through a
compromised Ukrainian accounting software update (supply-chain initial access),
then destroying backups and MBRs across hundreds of organisations worldwide.
The takeaway here is that a cyber weapon can cause catastrophic collateral
damage well beyond its intended geographic or organisational target: Maersk,
Merck, and FedEx were all affected despite not being the primary political target.
The CVSS~9.8 vulnerabilities studied in Week~4 (particularly the notion of
network-reachable, no-preconditions exploits) map directly to how NotPetya spread
laterally once inside a network perimeter.


\section{Lab Findings}

\subsection{Exercise 1: MITRE ATT\&CK Threat Actor Mapping}

Using the ATT\&CK Navigator, I built heat-map layers for Sandworm Team (APT44,
attributed to Russian GRU Unit 74455) and Lazarus Group (APT38, attributed to
North Korea's RGB). The structured data is encoded in \texttt{threat\_actor\_mapping.py}.

\vspace{4pt}
\noindent\textbf{Five Sandworm techniques selected:}

\small
\begin{tabularx}{\linewidth}{>{\ttfamily}l l >{\raggedright\arraybackslash}X}
\toprule
\normalfont ID & Tactic & Technique / observation \\
\midrule
T1486     & Impact (TA0040)             & Data Encrypted for Impact --- NotPetya/BadRabbit: encryption is cover for destruction, no recovery key exists \\
T1059.003 & Execution (TA0002)          & Windows Command Shell --- cmd.exe used to stage wipers and disable AV \\
T1078     & Defence Evasion (TA0005)    & Valid Accounts --- stolen credentials used for lateral movement, blends with normal traffic \\
T1071.001 & C2 (TA0011)                 & Web Protocols --- HTTPS C2 blends with legitimate browsing, hard to block without breaking access \\
T1561.002 & Impact (TA0040)             & Disk Structure Wipe --- Industroyer2 overwrites MBR; recovery requires full reinstall \\
\bottomrule
\end{tabularx}
\normalsize

\vspace{6pt}
\noindent\textbf{Comparison with Lazarus Group (3 overlapping techniques):}
Both groups use T1486 (wiper/ransomware), T1078 (valid accounts), and T1071
(application layer C2). This overlap is instructive rather than diagnostic: these
are among the most effective and widely documented techniques in ATT\&CK. Shared
TTPs do not imply shared origin. Attribution requires corroborating evidence:
infrastructure reuse, malware code similarity, timing patterns relative to
geopolitical events, and HUMINT. The overlap illustrates why attribution is
genuinely difficult and why confident public attribution is rare.

\subsection{Exercise 2: Disinformation Analysis}

I analysed Operation Secondary Infektion using the Stanford Internet Observatory /
EU DisinfoLab joint report as my primary source. The operation ran from at least
2014 to 2020, planted forged documents across 300+ platforms in 7 languages, and
targeted audiences in Ukraine, Germany, France, Spain, and the UK.

\noindent\textbf{Campaign anatomy:} Fabricated content included forged EU and NATO
official documents, fake news articles mimicking established outlets, and
manipulated images. Laundering occurred through newly registered typosquatting
domains. Amplification relied on coordinated cross-posting and circular citation
between fake sites to create an illusion of independent corroboration.

\noindent\textbf{OSINT tool applied:} WHOIS/DNS lookup (lookup.icann.org) on a
domain cited in the report. Key findings: (i) registration date preceded first
article by two weeks only, consistent with purpose-built infrastructure;
(ii) name server shared with two other fake outlets identified in the same
investigation --- the single most useful finding, linking three ``independent''
sites to the same actor. Limitation: WHOIS alone cannot prove intent;
privacy-shielded registration is legal and common.

\noindent\textbf{Red-flag indicator table (6 indicators):}

\small
\begin{tabularx}{\linewidth}{r >{\raggedright}p{3.8cm} >{\raggedright\arraybackslash}X l}
\toprule
\# & Indicator & How to check & Confidence \\
\midrule
1 & Coordinated posting timing & Botometer; cross-platform timestamp comparison & High \\
2 & Domain registration anomaly & WHOIS lookup; compare reg. date vs. first article & High \\
3 & Document metadata mismatch & exiftool / PDF properties; compare software locale & High \\
4 & Cross-platform self-citation loop & archive.org; trace content to single origin & Medium \\
5 & Linguistic inconsistency & Back-translation; native speaker review & Medium \\
6 & Infrastructure clustering & Reverse-IP lookup; certificate transparency logs & High \\
\bottomrule
\end{tabularx}
\normalsize

\subsection{Exercise 3: Attack Scenario and Countermeasures}

\noindent\textbf{Target:} Hospital / emergency medical services network.
\noindent\textbf{Actor:} Iron Caduceus (fictional nation-state APT, advanced,
state-level resources, politically motivated coercion).

The 7-step chain below maps to the Lockheed Martin Cyber Kill Chain. Full
structured output is in \texttt{attack\_scenario.py}.

\small
\begin{tabularx}{\linewidth}{r >{\raggedright}p{2.4cm} >{\raggedright}p{3.8cm} >{\raggedright\arraybackslash}X}
\toprule
Step & Stage & Attacker action & Countermeasure \\
\midrule
1 & Reconnaissance  & Passive OSINT: LinkedIn, job posts, conference slides & OSINT audit of own public footprint; social media policy \\
2* & Weaponisation   & Spear-phish impersonating EHR vendor; macro dropper & Disable macros by GPO; email sandboxing \\
3 & Delivery        & Email to IT admins and procurement; timed Monday AM & Phishing simulation; DMARC/DKIM/SPF \\
4 & Exploitation    & Macro spawns PowerShell; in-memory LOTL execution & EDR monitoring process lineage; PS script-block logging \\
5 & Installation    & Scheduled task + registry Run key; internal recon & SIEM alerts on task creation; network segmentation \\
6 & C\&C            & HTTPS C2, domain fronting, randomised beacon interval & TLS inspection; DNS threat intelligence feeds \\
7 & Actions         & Lateral movement to EHR; wiper destroys records and backups & Offline immutable backups (3-2-1); PAW for EHR admin \\
\bottomrule
\end{tabularx}
\normalsize

\vspace{4pt}
\noindent*\,\textbf{Weakest link: Step~2 (Weaponisation/Delivery).} A staff member
who recognises the phishing email and reports it before opening the attachment
breaks the entire chain before any code executes. At every subsequent step the
attacker already has a foothold, making defence progressively more expensive.
Investment in training and email filtering at this stage has the highest
defensive return on investment.


\section{Reflection and Critical Analysis}

Studying cybercrime --- as we did implicitly in Week~4 through CVE analysis ---
focuses on technical vulnerabilities: a specific version of vsftpd has a known
backdoor, and the fix is to update it. The exploitation is impersonal and largely
automated. Cyber terrorism and information warfare are different in a way that
matters for how a computer scientist thinks about defence.

The technical underpinning is often identical. NotPetya's lateral spread used
EternalBlue, the same SMB exploit class as the CVE-2007-2447 Samba vulnerability
examined in Week~4; the CVSS vector (network-reachable, no credentials required)
is the same category of risk. What changes is the adversary model: a ransomware
gang will stop if the victim pays or if the exploit is patched. An APT acting
on political orders does not stop. It adapts, waits, and comes back through a
different vector. This fundamentally changes the defensive posture required:
patching is still necessary, but it is not sufficient when the adversary has
the patience and resources to find the next unpatched system.

The disinformation exercise reinforced a different angle: the most effective
information operations are not technically sophisticated. Operation Secondary
Infektion did not require any CVEs. It required patience, language skills, and
an understanding of how human trust in media works. The OSINT indicators I
documented (metadata, registration timing, shared infrastructure) are technically
checkable by any analyst --- the defence against information warfare is partly
technical and partly a matter of critical reading skills that have nothing to
do with network security.

The weakest link in the hospital scenario (spear-phishing) connects directly
to the Principle of Least Privilege discussion in Week~2: the reason spear-phishing
is so effective against high-privilege accounts is precisely because those accounts
have been granted rights that far exceed daily task requirements. Tightening
access reduces the blast radius when the phish succeeds.

Looking ahead, information warfare intersects with industrial espionage in a
specific way: state-sponsored actors conducting IW campaigns often simultaneously
collect intellectual property and trade secrets from the same targets, using the
same infrastructure --- the two objectives are pursued in a single operation.


\section*{References}

\setlength{\parindent}{0pt}
\setlength{\parskip}{4pt}

\hangindent=1.27cm
Center for Internet Security. (2024). \textit{CIS Controls v8}. \url{https://www.cisecurity.org/controls}

\hangindent=1.27cm
Cybersecurity and Infrastructure Security Agency [CISA]. (2022). \textit{Alert AA22-057A: Destructive malware targeting organizations in Ukraine}. \url{https://www.cisa.gov/news-events/cybersecurity-advisories/aa22-057a}

\hangindent=1.27cm
European Union Agency for Cybersecurity [ENISA]. (2023). \textit{ENISA threat landscape 2023}. \url{https://www.enisa.europa.eu/topics/cyberthreats/enisa-threat-landscape}

\hangindent=1.27cm
Greenberg, A. (2019). \textit{Sandworm: A new era of cyberwar and the hunt for the Kremlin's most dangerous hackers}. Doubleday.

\hangindent=1.27cm
Lockheed Martin. (2022). \textit{Cyber kill chain}. \url{https://www.lockheedmartin.com/en-us/capabilities/cyber/cyber-kill-chain.html}

\hangindent=1.27cm
MITRE Corporation. (2024). \textit{ATT\&CK enterprise matrix v15}. \url{https://attack.mitre.org}

\hangindent=1.27cm
National Institute of Standards and Technology [NIST]. (2018). \textit{Framework for improving critical infrastructure cybersecurity} (Version 1.1). \url{https://www.nist.gov/cyberframework}

\hangindent=1.27cm
NATO Cooperative Cyber Defence Centre of Excellence [CCDCOE]. (2023). \textit{Tallinn manual 2.0 on the international law applicable to cyber operations}. \url{https://ccdcoe.org/research/tallinn-manual/}

\hangindent=1.27cm
Rid, T. (2020). \textit{Active measures: The secret history of disinformation and political warfare}. Farrar, Straus and Giroux.

\hangindent=1.27cm
Stanford Internet Observatory \& EU DisinfoLab. (2020). \textit{Secondary infektion: A large-scale cross-platform influence operation}. \url{https://secondaryinfektion.org}

\end{document}
